% GENERATED by generate_model.py from lsst/rubin-sizing-model @ v2026.09.2. Do not edit by hand.
% Regenerate: git clone --branch v2026.09.2 https://github.com/lsst/rubin-sizing-model ; then, inside the clone, uv run generate_model.py --emit-tex <this directory>

\small \begin{longtable}{|p{0.55\textwidth}|r|l|}
\caption{Observing and DRP compute inputs to the operations model, with the USDF batch fleet. \label{tab:smComputeInputs}}\\
\hline
\textbf{Metric} & \textbf{Value} & \textbf{Unit} \\
\hline\endfirsthead
\hline
\textbf{Metric} & \textbf{Value} & \textbf{Unit} \\
\hline\endhead
Observing nights per year & 300 & nights \\ \hline
Visits per night & 700 & visits \\ \hline
Image size (logical) & 0.00819222 & TB \\ \hline
Raw compression ratio & 0.464 & ratio \\ \hline
Input data per survey year & 798 & TB \\ \hline
DR1 compute estimate & 11,000 & node-days \\ \hline
Visits in that estimate & 188,340 & visits \\ \hline
One survey year vs that visit list & 1.115 & x \\ \hline
DR1 compute estimate in core-hours & 31,680,000 & core-hrs \\ \hline
DRP core-hours per input TB & 44,251 & ch/TB \\ \hline
DP2 compute (actual, on Milano) & 970 & node-days \\ \hline
Cores per batch node & 120 & cores \\ \hline
DRP processing window & 200 & days \\ \hline
Idle / inter-stage multiplier & 1.35 & x \\ \hline
USDF RSP extra load fraction & 1.00\% & fraction \\ \hline
Developer / pilot batch load & 55 & node-days/yr \\ \hline
Safety margin & 1.5 & x \\ \hline
Batch nodes & 145 & nodes \\ \hline
Milano nodes & 110 & nodes \\ \hline
Torino nodes & 35 & nodes \\ \hline
Torino per-core speedup vs Milano & 2.0 & x \\ \hline
Fleet capacity, Milano-equivalent cores & 21,600 & cores \\ \hline
\end{longtable} \normalsize
